\chapter{三角函数}\label{ch:g12:trigo}

函数 $\cos$ 与 $\sin$ 把圆的几何
转化为分析问题。它们通过将实数轴绕在单位圆上而定义，
是周期函数的典型，其导数使它们
成为振动方程 $y'' = -y$ 的解。

\section{单位圆}

\begin{definition}[余弦与正弦]\label{def:g12:trigo:cossin}
设 $\mathcal{C}$ 为标准正交坐标系中以原点为圆心、
半径为 1 的圆。对每个实数 $t$，取点 $M(t)$，
取长为 $\abs{t}$ 的弧，将其绕在 $\mathcal{C}$ 上，
起点为 $I(1,0)$，当 $t \geq 0$ 时按逆时针方向，否则按顺时针方向。则
\emph{$\cos t$ 与 $\sin t$ 是点
$M(t)$ 的坐标}\index{余弦}\index{正弦}：
\[
M(t) = (\cos t, \sin t).
\]
\end{definition}

\begin{omfigure}
\begin{tikzpicture}[scale=2.2]
  \draw[->] (-1.25,0) -- (1.35,0) node[right] {$x$};
  \draw[->] (0,-1.25) -- (0,1.35) node[above] {$y$};
  \draw (0,0) circle (1);
  \coordinate (O) at (0,0);
  \coordinate (M) at (50:1);
  \draw[thick,omDef] (O) -- (M);
  \fill (M) circle (0.6pt) node[above right] {$M(t)$};
  \fill (1,0) circle (0.6pt) node[below right] {$I$};
  \draw[dashed] (M) -- (M |- O) node[below,pos=1.05] {$\cos t$};
  \draw[dashed] (M) -- (M -| O) node[left,pos=1.02] {$\sin t$};
  \draw[->,omThm,thick] (0.35,0) arc (0:50:0.35);
  \node[omThm] at (25:0.5) {$t$};
\end{tikzpicture}
\end{omfigure}

由于圆的周长为 $2\pi$，点 $M(t)$
在 $t$ 增加 $2\pi$ 时保持不变；又因 $M(t)$ 位于单位圆上，有

\begin{proposition}[基本恒等式]\label{prop:g12:trigo:identities}
对所有 $t \in \R$ 与 $k \in \Z$，有
\[
\cos^2 t + \sin^2 t = 1, \qquad
\cos(t + 2k\pi) = \cos t, \qquad
\sin(t + 2k\pi) = \sin t,
\]
\[
\cos(-t) = \cos t, \qquad \sin(-t) = -\sin t,
\]
\[
\cos(\pi - t) = -\cos t, \quad \sin(\pi - t) = \sin t, \quad
\cos\left(\tfrac{\pi}{2} - t\right) = \sin t, \quad
\sin\left(\tfrac{\pi}{2} - t\right) = \cos t .
\]
\end{proposition}

\begin{proof}
第一个恒等式即单位圆的方程；其余各式表达圆的对称性：
$M(-t)$ 是 $M(t)$ 关于
$x$ 轴的对称点，$M(\pi - t)$ 关于 $y$ 轴的对称点，$M(\frac\pi2 - t)$ 关于
直线 $y = x$ 的对称点。
\end{proof}

需要熟记的数值：
\begin{center}
\begin{tabular}{c|ccccc}
$t$ & $0$ & $\dfrac{\pi}{6}$ & $\dfrac{\pi}{4}$ & $\dfrac{\pi}{3}$ & $\dfrac{\pi}{2}$\\[6pt]
\hline\rule{0pt}{14pt}%
$\cos t$ & $1$ & $\dfrac{\sqrt3}{2}$ & $\dfrac{\sqrt2}{2}$ & $\dfrac12$ & $0$\\[6pt]
$\sin t$ & $0$ & $\dfrac12$ & $\dfrac{\sqrt2}{2}$ & $\dfrac{\sqrt3}{2}$ & $1$
\end{tabular}
\end{center}

\begin{proposition}[加法公式]\label{prop:g12:trigo:addition}
对所有 $a, b \in \R$，有
\begin{align*}
\cos(a+b) &= \cos a \cos b - \sin a \sin b, &
\sin(a+b) &= \sin a \cos b + \cos a \sin b,\\
\cos(a-b) &= \cos a \cos b + \sin a \sin b, &
\sin(a-b) &= \sin a \cos b - \cos a \sin b.
\end{align*}
特别地，$\cos 2a = 2\cos^2 a - 1 = 1 - 2\sin^2 a$ 且
$\sin 2a = 2 \sin a \cos a$。
\end{proposition}

\begin{proof}
向量 $\vect{u} = (\cos a, \sin a)$ 与 $\vect{v} = (\cos b, \sin b)$ 为单位向量，
它们的夹角为 $a - b$；其数量积一方面等于
$\cos a\cos b + \sin a \sin b$（用坐标计算），另一方面等于
$\norm{\vect u}\,\norm{\vect v}\cos(a-b) = \cos(a - b)$（用几何方法计算），这就是第三个公式。
将 $b$ 换为 $-b$ 即得第一个公式；正弦公式
利用 $\sin x = \cos(\frac\pi2 - x)$ 推出。倍角公式是取
$b = a$ 并结合 $\cos^2 + \sin^2 = 1$ 所得。
\end{proof}

\section{分析研究}

\begin{lemma}[基本极限]\label{lem:g12:trigo:sinxx}
\[
\lim_{t \to 0} \frac{\sin t}{t} = 1
\qquad\text{且}\qquad
\lim_{t \to 0} \frac{\cos t - 1}{t} = 0 .
\]
\end{lemma}

\begin{proof}
对 $0 < t < \frac{\pi}{2}$，比较单位圆中的三块面积：
三角形 $OIM(t)$、扇形 $OIM(t)$，以及
底边为 $OI$、高为 $\tan t$ 的直角三角形：
\[
\frac{\sin t}{2} \leq \frac{t}{2} \leq \frac{\tan t}{2}.
\]

\begin{omfigure}
\begin{tikzpicture}[scale=2.2]
  \draw[->] (-0.25,0) -- (1.5,0) node[right] {$x$};
  \draw[->] (0,-0.25) -- (0,1.45) node[above] {$y$};
  \fill[omDef!20] (0,0) -- (1,0) -- (50:1) -- cycle;
  \draw[gray, thin] (-10:1) arc (-10:80:1);
  \draw[omProp, thick] (0,0) -- (1,0) -- (1,1.1918) -- cycle;
  \draw[omDef] (0,0) -- (1,0) -- (50:1) -- cycle;
  \draw[omThm, thick] (1,0) arc (0:50:1);
  \draw[dashed] (50:1) -- (0.643,0);
  \node[left, font=\small] at (0.68,0.35) {$\sin t$};
  \node[omProp, right, font=\small] at (1,0.62) {$\tan t$};
  \node[omThm, font=\small] at (20:0.5) {$t$};
  \fill (50:1) circle (0.6pt);
  \node at (58:1.22) {$M(t)$};
  \fill (1,0) circle (0.6pt) node[below right] {$I$};
  \node[below left] at (0,0) {$O$};
\end{tikzpicture}

{\small 三个嵌套的区域：三角形 $OIM$（蓝色，面积
$\frac{\sin t}{2}$）、扇形（红色弧线所界，面积
$\frac{t}{2}$）以及高为 $\tan t$ 的直角三角形（橙色，面积
$\frac{\tan t}{2}$）。}
\end{omfigure}

由第一个不等式得 $\frac{\sin t}{t} \leq 1$；由第二个得
$\cos t \leq \frac{\sin t}{t}$。当 $t \to 0^+$ 时，由连续性 $\cos t \to 1$，
由夹逼定理得 $\frac{\sin t}{t} \to 1$，利用奇偶性，
可得 $t \to 0$ 时的结论。对第二个极限，
\[
\frac{\cos t - 1}{t} = \frac{\cos^2 t - 1}{t(\cos t + 1)}
= -\frac{\sin t}{t}\cdot\frac{\sin t}{\cos t + 1}
\xrightarrow[t\to0]{} -1 \cdot \frac{0}{2} = 0. \qedhere
\]
\end{proof}

\begin{theorem}[正弦与余弦的导数]\label{thm:g12:trigo:derivatives}
函数 $\sin$ 与 $\cos$ 在 $\R$ 上可导，且
\[
(\sin)' = \cos, \qquad (\cos)' = -\sin .
\]
一般地，$(\sin(ax+b))' = a\cos(ax+b)$，
$(\cos(ax+b))' = -a\sin(ax+b)$。
\end{theorem}

\begin{proof}
由和角公式，
\[
\frac{\sin(x+h) - \sin x}{h}
= \sin x\,\frac{\cos h - 1}{h} + \cos x\,\frac{\sin h}{h}
\xrightarrow[h\to0]{} \sin x \cdot 0 + \cos x \cdot 1 = \cos x,
\]
利用\cref{lem:g12:trigo:sinxx}。$\cos$ 的计算完全类似，
一般形式由链式法则得到。
\end{proof}

\begin{proposition}[单调性]\label{prop:g12:trigo:variations}
在 $\intcc{0}{\pi}$ 上，$\cos$ 从 $1$ 递减到 $-1$；
在 $\intcc{-\frac\pi2}{\frac\pi2}$ 上，$\sin$ 从 $-1$ 递增到 $1$。两个
函数都是 $2\pi$ 周期的，值域均为 $\intcc{-1}{1}$；$\cos$ 是
偶函数，$\sin$ 是奇函数。
\end{proposition}

\begin{proof}
在 $\intoo{0}{\pi}$ 上，$(\cos)' = -\sin < 0$，因为点 $M(t)$ 在此处
的纵坐标为正；同理 $(\sin)' = \cos > 0$ 在
$\intoo{-\frac\pi2}{\frac\pi2}$ 上成立。周期性与奇偶性由
\cref{prop:g12:trigo:identities}可得。
\end{proof}

\begin{omfigure}
\begin{tikzpicture}
\begin{axis}[
  width=12cm, height=4.5cm,
  axis lines=middle,
  xmin=-7, xmax=7, ymin=-1.4, ymax=1.4,
  xtick={-6.2832,-3.1416,3.1416,6.2832},
  xticklabels={$-2\pi$,$-\pi$,$\pi$,$2\pi$},
  ytick={-1,1},
  samples=200, domain=-7:7]
  \addplot[omDef,thick] {cos(deg(x))};
  \addplot[omThm,thick] {sin(deg(x))};
\end{axis}
\end{tikzpicture}

{\small $\cos$ 的图像（蓝色）与 $\sin$ 的图像（红色）。}
\end{omfigure}

\begin{method}[解方程 $\cos x = a$ 与 $\sin x = a$]\label{met:g12:trigo:equations}
对 $a \in \intcc{-1}{1}$，求出一个特解 $\alpha$（可由三角函数表或
计算器得到）。则全部解为：
\[
\cos x = \cos\alpha \iff x = \alpha + 2k\pi \text{ 或 } x = -\alpha + 2k\pi,
\quad k \in \Z;
\]
\[
\sin x = \sin\alpha \iff x = \alpha + 2k\pi \text{ 或 } x = \pi - \alpha + 2k\pi,
\quad k \in \Z.
\]
解不等式时，在单位圆上找出解所在的弧，
并在一个周期内读出相应区间。
\end{method}

\begin{example}\label{ex:g12:trigo:equation}
求方程 $\cos x = \frac12$ 在 $\intoc{-\pi}{\pi}$ 内的解：
一个特解为 $\frac\pi3$，故解为 $x = \frac\pi3$ 与 $x = -\frac\pi3$。
在 $\R$ 上全部解为：$x = \pm\frac\pi3 + 2k\pi$，$k \in \Z$。
\end{example}

\section{练习}

\begin{exercise}[$\star$]\label{exo:g12:trigo:1}
计算 $\cos\frac{2\pi}{3}$、$\sin\frac{5\pi}{6}$、$\cos\frac{7\pi}{4}$ 及
$\sin\left(-\frac{\pi}{3}\right)$ 的值，
利用\cref{prop:g12:trigo:identities}中的公式。
\end{exercise}

\begin{exercise}[$\star$]\label{exo:g12:trigo:2}
求下列函数的导数：$f(x) = \sin^2 x$、$g(x) = \cos(3x) + x\sin x$ 及
$h(x) = \dfrac{\sin x}{2 + \cos x}$。
\end{exercise}

\begin{exercise}[$\star$]\label{exo:g12:trigo:3}
在 $\intoc{-\pi}{\pi}$ 内、再在 $\R$ 上求下列方程的解：
\[
\text{(a) } \sin x = \frac{\sqrt3}{2};
\qquad
\text{(b) } \cos\left(2x\right) = \frac{\sqrt2}{2} .
\]
\end{exercise}

\begin{exercise}[$\star\star$]\label{exo:g12:trigo:4}
利用和角公式，求
$\cos\frac{\pi}{12}$ 与 $\sin\frac{\pi}{12}$ 的精确值。
（提示：$\frac{\pi}{12} = \frac{\pi}{3} - \frac{\pi}{4}$。）
\end{exercise}

\begin{exercise}[$\star\star$]\label{exo:g12:trigo:5}
在 $\intcc{0}{2\pi}$ 内解不等式 $2\sin^2 x - \sin x - 1 \geq 0$。
（提示：把关于 $\sin x$ 的二次多项式因式分解。）
\end{exercise}

\begin{exercise}[$\star\star$]\label{exo:g12:trigo:6}
求下列极限：
\[
\lim_{x\to0} \frac{\sin 3x}{x}, \qquad
\lim_{x\to0} \frac{1 - \cos x}{x^2}, \qquad
\lim_{x\to+\infty} x \sin\frac{1}{x}.
\]
\end{exercise}

\begin{exercise}[$\star\star$]\label{exo:g12:trigo:7}
研究函数 $f(x) = x + \cos x$ 在 $\intcc{0}{2\pi}$ 上的单调性与
极值；证明 $f$ 在 $\R$ 上递增，尽管 $f'$ 在
无穷多个点处为零。
\end{exercise}

\begin{exercise}[$\star\star\star$]\label{exo:g12:trigo:8}
设 $f(x) = A\cos x + B \sin x$，其中 $(A,B) \neq (0,0)$。
\begin{enumerate}
  \item 证明 $f$ 可写成 $f(x) = R\cos(x - \varphi)$ 的形式，其中
        $R = \sqrt{A^2 + B^2}$，且 $\varphi$ 为适当选取的值。
  \item 由此求 $f$ 的最大值与最小值，并
        解方程 $\cos x + \sin x = 1$（在 $\intoc{-\pi}{\pi}$ 内）。
\end{enumerate}
\end{exercise}

\begin{exercise}[$\star\star\star$]\label{exo:g12:trigo:9}
证明 $\sin$ 与 $\cos$ 都满足微分方程
$y'' = -y$。反之，设 $y$ 为 $y'' = -y$ 的解，满足
$y(0) = a$、$y'(0) = b$，证明 $y = a\cos + b\sin$。
（提示：考虑 $g = (y - a\cos - b\sin)$ 及“能量”
$E = g^2 + (g')^2$。）
\end{exercise}

\section{问题：摆的秘密}

\begin{problem}[{周末问题 --- 一个几何极限支撑了
正弦与余弦的导数，并由此导出整个振动：
弹簧、摆、拍，以及几乎
成为“米”的那个摆长}]
\label{pb:g12:trigo:1}
凡曾滴答作响的钟，凡曾发声的弦，
都服从同一套数学：回复力与
位移成正比，于是有方程 $y'' = -\omega^2 y$，进而得出
余弦。其根基是一个看似平凡的极限，
$\frac{\sin x}{x} \to 1$ --- 九年级时已预告，并在
\cref{thm:g12:trigo:derivatives}中用到，本题借助图形予以证明。
本题由该极限出发，直至一米长摆的摆动，
并说明为何“米”差之毫厘，
未能由此摆来定义。

\medskip
\noindent\textbf{第一部分 --- 基本极限。}
\begin{enumerate}
  \item 先做数值热身：计算 $\frac{\sin x}{x}$ 在
        $x = 0.1$ 与 $x = 0.01$ 处的值（弧度制！）。
  \item 夹逼论证：在单位圆上，
        取 $0 < x < \frac\pi2$，用初等几何表示下列图形的面积：
        （i）顶点为 $O$、$A(1,0)$
        及圆上点 $M(\cos x, \sin x)$ 的三角形；（ii）
        扇形 $OAM$；（iii）直角三角形 $OAT$，
        其中 $T(1, \tan x)$。
  \item 由面积的夹逼关系可得
        $\sin x < x < \tan x$，则
        $\cos x < \frac{\sin x}{x} < 1$，由夹逼定理
        可知
        $\lim_{x \to 0} \frac{\sin x}{x} = 1$（奇偶性处理
        $x < 0$ 的情形）。
  \item 由此可得
        $\lim_{x \to 0} \frac{1 - \cos x}{x^2} = \frac12$
        （乘以共轭式）。
  \item 说明为何这一极限是
        \cref{thm:g12:trigo:derivatives}的关键（
        $\sin'(0)$ 为多少？），并说明为何它还
        解决了一个老问题：在何种意义上，\emph{弧度}是唯一使
        $\sin' = \cos$ 无需常数因子即成立的角度单位？
\end{enumerate}

\medskip
\noindent\textbf{第二部分 --- 简谐运动。}
\begin{enumerate}[resume]
  \item 用链式法则验证
        $y(t) = R\cos(\omega t - \varphi)$ 满足
        $y'' = -\omega^2 y$。
  \item 一个挂在弹簧上的质点满足 $y'' = -4y$，且
        $y(0) = 3$，$y'(0) = 8$。由
        \cref{exo:g12:trigo:9}可承认所有解均为
        $A\cos 2t + B\sin 2t$，并求 $A$ 与 $B$。
  \item 将解改写为 $R\cos(2t - \varphi)$
        （见\cref{exo:g12:trigo:8}），给出振幅 $R$、
        周期，以及相位 $\varphi$（保留三位小数）。
  \item 质点于何时首次经过
        平衡位置 $y = 0$？
  \item 摆：长为 $L$ 的细绳系一摆锤，它
        \emph{精确地}满足 $\theta'' = -\frac gL \sin\theta$ ---
        用初等函数无法求解。对于小幅摆动，
        将 $\sin\theta$ 换为 $\theta$（第一部分已对此作出论证！），
        并由此得出周期
        $T = 2\pi\sqrt{\frac Lg}$，并求 $T$ 在
        $L = 1$ m（$g = 9.81$）时的值：伽利略的著名发现 ---
        公式的结果\emph{不}依赖于什么？
  \item 秒摆：摆长取何值时
        恰好有 $T = 2$ s？1791 年，法国科学院曾在
        此摆长与四分之一子午线的千万分之一之间犹豫，
        以定义\emph{米}——
        两个候选值相差多少毫米？
\end{enumerate}

\medskip
\noindent\textbf{第三部分 --- 和角公式的应用。}
\begin{enumerate}[resume]
  \item 利用和角公式
        （\cref{prop:g12:trigo:addition}）证明
        $\cos\left(x + \frac\pi2\right) = -\sin x$ 与
        $\sin\left(x + \frac\pi2\right) = \cos x$：三角求导相当于
        转过四分之一圈——请检验其与
        \cref{thm:g12:trigo:derivatives}是否一致。
  \item 推导积化和差公式
        $\cos a \cos b = \frac12\left(\cos(a - b) +
        \cos(a + b)\right)$，以及降幂公式
        $\cos^2 x = \frac{1 + \cos 2x}{2}$（积分一章
        会感谢你的）。
  \item 拍：推导
        $\sin p + \sin q = 2 \sin\frac{p + q}{2}
        \cos\frac{p - q}{2}$，并将其用于两个音叉
        频率为 $440$ 与 $444$ Hz：耳朵听到的是什么音，
        每秒有多少次拍？（这正是乐团调音所用的
        方法。）
  \item 推导三倍角公式
        $\cos 3x = 4\cos^3 x - 3\cos x$
        （写作 $3x = 2x + x$）。令 $x = 20^\circ$，
        $\cos 20^\circ$ 应满足怎样的三次方程？（这个
        三次方程仅用平方根无法求解，这正是
        角的三等分使尺规作图束手无策两千年
        ——完整的故事
        见大学分册。）
  \item 某交流电的电流为
        $I(t) = 3\cos(100\pi t) + 3\sqrt3 \sin(100\pi t)$.
        将其改写为 $R\cos(100\pi t - \varphi)$，求
        峰值电流与相位？
\end{enumerate}

\medskip
\noindent\textbf{第四部分 --- 振动的世界。}
\begin{enumerate}[resume]
  \item 求下列函数的（最小）周期：
        $\sin^2 t$（利用第 13 题）；
        $\sin(2t) + \sin(3t)$ 呢？
  \item 一个真实的摆存在阻尼：
        $y(t) = \eu^{-t/10}\cos(2\pi t)$。描述其运动，
        计算 $t = 10$ 时包络线的值（一位老朋友
        再次出现，见\cref{pb:g12:exp:1}），并指出
        此类解所对应的方程在哪一章中求解。
  \item 为什么正弦与余弦支配着一切振动？用两句话给出
        从物理到数学的对照表，
        （回复力与位移成正比
        $\rightarrow$ 对应哪个方程 $\rightarrow$
        哪些解），并说明\cref{exo:g12:trigo:9}
        对这一论断有何贡献。
  \item 总结：本章的脉络分四步——一个几何
        极限，两次求导，一个微分方程，
        以及整个振动问题。诚实的附注：
        对于大幅摆动，摆的真实周期涉及一个
        椭圆积分，借助\,\dots 以闪电般的速度算出：
        高斯的算术--几何平均数，即
        \cref{pb:g12:seq:1}中的甜点。各条线索自行连接。
\end{enumerate}
\end{problem}
